\section{题03：多项式积分（P16785）}

\subsection{问题重述}

题目给出一行字符串表示一个关于 $x$ 的一元多项式，要求输出其不定积分（不含积分常数 $C$）。输入多项式各项按次数从高到低排列，首项为正时前面不带 $+$ 号，首项为负时带 $-$ 号，后续各项以 $+$ 或 $-$ 分隔且符号归属于该项。系数为 $\pm 1$ 时可省略数字 $1$（例如 $x^{\wedge}3$ 表示 $1x^3$，$-x$ 表示 $-1x$），次数为 $1$ 时写作 $cx$ 或 $x$，次数为 $0$ 时写作整数常数，次数 $\ge 2$ 时写作 $cx^{\wedge}k$ 或 $x^{\wedge}k$。输出也需遵守类似格式，且系数为分数 $\frac{p}{q}$ 时要输出为 $p/q$，要求 $p,q$ 互质、$q>0$，负号放在分子上；若化简后分母为 $1$ 则直接输出整数。约束：项数 $\le 10^4$，次数 $\le 10^8$，系数绝对值 $\le 10^8$。

以样例 $1$ 为例，输入 \mil{4x\^{}5-4x\^{}3+2}，不定积分为 $\int(4x^5-4x^3+2)\,dx = \frac{4}{6}x^6 + \frac{-4}{4}x^4 + \frac{2}{1}x^1 = \frac{2}{3}x^6 - x^4 + 2x$，输出 \mil{2/3x\^{}6-x\^{}4+2x}。这个例子已经展示了问题的全部要素：系数从整数变成既约分数、分母为 $1$ 时省略分母、系数为 $-1$ 时省略数字 $1$。

\subsection{积分的线性性：逐项独立处理的数学依据}

不定积分的幂法则告诉我们 $\int cx^k \,dx = \frac{c}{k+1}x^{k+1} + C$。不定积分是一个线性算子，即 $\int (f+g) = \int f + \int g$ 且 $\int \alpha f = \alpha \int f$，因此多项式积分可以对每一项独立计算后直接拼接结果，各项之间互不干扰。更重要的是，每一项的次数加 $1$ 后，原来的降序排列自然保持——$k_1 > k_2$ 蕴含 $k_1+1 > k_2+1$，因此不需要额外排序。

于是整个问题分解为三个独立的阶段：将输入字符串解析为若干 $(c_i, k_i)$ 对；对每一对计算 $p_i = c_i,\; q_i = k_i+1,\; d_i = k_i+1$ 并用 $\gcd(|p_i|, q_i)$ 约简；按输出格式规则将结果拼接为字符串。数学核心只占三行代码，真正需要细致处理的是解析和格式化的诸多边界约定。

\subsection{解析的暗礁：隐式系数、符号归属与三种次数表示}

解析是本题最容易出错的部分。输入格式有若干“隐式约定”，需要逐字符扫描时精确判断。

符号归属规则：首项若为正则前面没有符号字符，首项若为负则有 $-$ 号；非首项前面一定有 $+$ 或 $-$。因此解析循环的每一轮，若当前字符是 $+$ 或 $-$ 则记录符号并前进指针，然后读取该项的系数和次数。

系数的隐式约定：若 $x$ 前面没有出现数字，则系数默认为 $1$（乘以符号后可能是 $-1$）。例如输入 \mil{x-1}，第一项在 \mil{x} 之前没有数字，系数为 $1$；第二项 \mil{-1} 有数字 $1$，系数为 $-1$。解析时用一个标志 \mil{have_coeff} 记录是否读到数字：读到了就用读到的值，没读到就默认为 $1$。读数字时逐位累加 \mil{coeff = coeff * 10 + digit}。

次数的三种表示：若当前字符是 \mil{x}，说明该项含有变量；接着若下一个字符是 \mil{\^{}} 则后面跟有指数数字（次数 $\ge 2$），否则次数为 $1$；若根本没遇到 \mil{x}，说明该项是常数，次数为 $0$。例如输入 \mil{4x\^{}5-4x\^{}3+2}，第一项的 \mil{x} 后跟 \mil{\^{}5}，次数为 $5$；最后一项 \mil{+2} 没有 \mil{x}，次数为 $0$。而 \mil{x} 单独出现时（如样例 $4$ 的 \mil{x-1}），第一项遇到 \mil{x} 且后面没有 \mil{\^{}}，次数为 $1$。

解析过程的正确性可以用几个边界输入检验：常数多项式 \mil{7} 应解析为 $(7,0)$；纯 $x$ 项 \mil{x} 应解析为 $(1,1)$；负系数无数字 \mil{-x} 应解析为 $(-1,1)$；大系数大次数如 \mil{-100x\^{}100} 应解析为 $(-100,100)$。解析结果的 \mil{vector<pair<ll,ll>>} 按输入顺序保存，即已经按次数降序排列。

\subsection{分数化简：符号归一化与分母恒正的约定}

对每一项 $(c_i, k_i)$，积分后的分子为 $c_i$，分母为 $k_i+1$，次数为 $k_i+1$。由于输入保证 $k_i \ge 0$，分母至少为 $1$。

约分步骤：设 $g = \gcd(|c_i|, k_i+1)$，用欧几里得算法在 $\mathrm{O}(\log \min(|c_i|, k_i+1))$ 时间内求出，然后分子分母同除以 $g$。注意 $\gcd$ 的参数应取绝对值，因为 $c_i$ 可能为负。

符号归一化步骤：输出规则要求分母恒为正，负号放在分子上。约分后若分母为负，则将分子分母同时取反。实现中只需判断 \mil{if (den < 0) \{ num = -num; den = -den; \}}。

以样例 $2$ 的第一项 $-100x^{100}$ 为例：$c = -100$，$k = 100$，分母 $101$，$\gcd(100, 101) = 1$（$100$ 和 $101$ 互质），约分后分子 $-100$，分母 $101$，次数 $101$。第二项 $105x^2$：$c=105$，$k=2$，分母 $3$，$\gcd(105,3)=3$，约分后分子 $35$，分母 $1$，次数 $3$——分母为 $1$ 表示系数化简为整数。两项组合输出为 \mil{-100/101x\^{}101+35x\^{}3}。

\subsection{输出格式的考验：首项符号、系数省略与分数表示}

输出阶段需要逐项判断三项内容：该项前面的符号、系数部分的字符串、变量部分的字符串。

符号规则分情况：若为第一项，正系数前面不加 $+$，负系数前面加 $-$；若为非首项，正系数加 $+$，负系数加 $-$。实现中只需检查 \mil{num} 的正负即可。

系数部分需要决定是否省略数字 $1$。规则是“单项系数为 $\pm 1$ 时省略数字 $1$”——这仅当分母为 $1$ 且分子绝对值为 $1$ 时才触发。例如积分结果中可能出现 $1x^3$，应输出为 \mil{x\^{}3}；$-1x$ 应输出为 \mil{-x}。对于分数系数（分母 $>1$），即使分子绝对值为 $1$，也不能省略，必须完整输出分子和分母。因为 $1/2$ 不是 $\pm 1$，省略 $1$ 会变成 \mil{/2}，无意义。因此代码逻辑为：若 \mil{den == 1} 且 \mil{abs_num == 1}，省略数字；若 \mil{den == 1} 且 \mil{abs_num != 1}，输出 \mil{abs_num}；若 \mil{den > 1}，输出 \mil{abs_num/den}。

变量部分：次数为 $1$ 时仅输出 \mil{x}（不输出指数 $^{\wedge}1$）；次数 $\neq 1$ 时输出 \mil{x\^{}deg}。注意常数项（次数 $0$）经过积分后次数变为 $1$，因此积分结果中不存在次数为 $0$ 的项——常数项的积分变成了 $x$ 的一次项，这正是样例 $3$ 中 $7$ 积分得到 $7x$ 的原因。

\subsection{复杂度分析与实现细节}

设输入字符串长度为 $n$，项数为 $m$（$m \le 10^4$）。每个字符最多被扫描一次，解析阶段 $\mathrm{O}(n)$。积分阶段每项做一次 $\gcd$，欧几里得算法复杂度 $\mathrm{O}(\log \max(|c_i|, k_i))$，由于 $\max(|c_i|, k_i) \le 10^8$，单次 $\gcd$ 约 $27$ 次迭代，总计 $m \times \log 10^8 \approx 2.7 \times 10^5$ 次运算。输出阶段 $\mathrm{O}(m)$。总体时间复杂度 $\mathrm{O}(n + m \log M)$，在给定约束下完全可行。

空间复杂度 $\mathrm{O}(m)$，存储解析结果和积分结果两组 \mil{vector}，每组约 $10^4$ 个元素，每个元素含三个 \mil{long long}，总计不足 $1$ MB。

实现中几个细节值得注意。使用 \mil{long long}（$64$ 位有符号整数）而非 \mil{int}，因为系数可达 $10^8$、次数可达 $10^8$，在计算 \mil{den = deg + 1} 和 \mil{gcd(num, den)} 时不会溢出，但若涉及乘法（如输出前做格式拼接）时 $64$ 位也足够。$\gcd$ 函数应先对参数取绝对值，因为分子可能为负，而欧几里得算法在负数上的行为是实现定义的。解析函数中标志 \mil{have_coeff} 至关重要：它区分了“系数为 $1$（省略数字）”和“常数的系数恰好是 $1$”这两种情形——对于输入 \mil{x}，没有数字但系数应为 $1$；对于输入 \mil{1}，有数字且系数就是 $1$。两种情形解析结果一致，但逻辑路径不同。最后，输出循环中注意首项与非首项的符号处理差异，以及分子绝对值 \mil{abs_num} 在符号已输出的前提下应使用绝对值。

\subsection{代码实现}
\inputminted{c++}{../src/p03_polynomial_integral.cpp}
